\section{把数学变成 Lean：形式化证明闭环}

把数学变成 Lean，不是简单翻译。自然语言证明常常省略背景、跳步、依赖图形直觉或引用不明确的定理。Lean 需要明确对象类型、定理陈述、依赖、tactic 和每一步合法性。这让 AI 既有机会，也有约束。

\begin{figure}[H]
\centering
\includegraphics[width=0.92\textwidth]{lean-proof-loop.png}
\caption{从自然语言证明到 Lean 验证的闭环。}
\end{figure}

\begin{knowledgebox}{读图：Lean 闭环为什么适合训练}
自然语言证明先被转成 formal statement，再生成 tactic 或 proof term，Lean kernel 检查是否合法。如果失败，系统返回 error/goal state，人或模型据此修复。这个闭环提供了比普通文本生成更强的训练反馈。
\end{knowledgebox}

\subsection{形式化的三类难点}

第一是语义难点：自然语言里的“显然”“由此可得”需要展开。第二是库依赖难点：Lean 里定理、定义、lemma 的名字和使用方式必须准确。第三是搜索难点：证明空间巨大，模型需要决定先证哪个子目标、调用哪个 lemma、何时重写目标。

\begin{warningbox}{Lean 不会自动让模型懂数学}
Lean 能验证证明是否合法，但不会替模型选择重要问题，也不会自动提供数学审美。AI for Math 还需要人类数学家的直觉、问题选择和理论品味。
\end{warningbox}

\subsection{数学天书中的证明}

访谈标题提到“数学天书中的证明”，可以理解为数学中那些优雅、短小、仿佛来自天书的证明。AI 若只是暴力搜索，可能产出冗长证明；若能学习数学直觉和审美，才可能接近这类证明。形式化验证保证正确性，但优雅性仍需要审美评价。

\begin{importantbox}{正确性与优雅性}
机器验证解决“对不对”，但不直接解决“好不好”。数学研究还关心证明是否解释了现象、是否揭示结构、是否可推广。AI for Math 的长期目标应同时追求正确性、可读性和洞察力。
\end{importantbox}

\subsection{本章小结}

Lean 让证明变成可执行、可检查对象。AI for Math 的关键，是把自然语言直觉和机器验证闭环连接起来，同时不丢掉数学审美。

